\chapter{A-TP via de Branges Hermite--Biehler Factorisation}
\label{v4-arith:w7-c:chapter}

\emph{The restricted-class comparison A (\Cref{v4-arith:w5-a:chapter}) reduced HP--Li on the
holomorphic cuspidal restricted class to the archimedean Tate
positivity conjecture A-TP, and
\Cref{v4-arith:w5-a:thm:ATP-equals-weil-positivity} placed A-TP
strictly equivalent to RH. The de Branges spectral comparison
(\Cref{v4-arith:w5-j:VBstar-spectral-positivity}) analysed the
main-direction arrow (HB\(^{+}\)) + (BD\(^{+}\)) + (Det) \(\Rightarrow\) RH
through the Hamburger--Carleman--de Branges spectral hypotheses. The bounded-interval reduction
(\Cref{v4-arith:w6-a:ATP-direct-obstruction}) analysed A-TP via
the digamma kernel \(g\), closed an explicit \(W_{\infty}\)-positive
HFD sub-class, and confined the sign-changing archimedean term to
\([-t_{\ast},t_{\ast}]\) with the arithmetic functional \(R[f]\)
retained. The de Branges Hermite--Biehler reformulation of A-TP,
the third non-elementary construction in
\Cref{v4-arith:w6-a:rem:live-constructions}, attaches to the completed
zeta the Lagarias function \(E_{\xi}=\xi_{\mathbb R}(\tfrac12-iz)+
\xi_{\mathbb R}'(\tfrac12-iz)\), identifies
\((E_{\xi}+E_{\xi}^{*})/2\) with \(\Xi\), places the Mellin transform
in the de Branges model space under HB\(^{+}\), translates the
Vinogradov--Korobov region under \(s=\tfrac12+iz\), and removes the
false central-strip closure.
The translation gives only peripheral zero-free bands in the
\(z\)-plane. It gives no positivity on the critical-line strip.
In this de Branges chart the internal positivity condition is
HB\(^{+}\), an RH-equivalent condition. The bounded-interval
archimedean term of the bounded-interval reduction becomes a Gram-positivity
problem for the de Branges kernel on the Mellin image; this gives the
isometry
\(\mathrm{A\text{-}TP}|_{\mathrm{PW}^{\mathrm{hol}}_{F}}\leftrightarrow
K_{E_{\xi}}|_{\mathrm{Mellin}(\mathrm{PW}^{\mathrm{hol}}_{F})}\) in
the arithmetic branch.}

\section{Context and the de Branges Comparison}
\label{v4-arith:w7-c:context}

Chapter~\ref{v4-arith:w5-a:chapter} established:
\begin{itemize}
\item A-TP \(\Leftrightarrow\) RH via the Weil positive cone
(\Cref{v4-arith:w5-a:thm:ATP-equals-weil-positivity}).
\item Prime-side Petersson closure on
\(\mathrm{PW}^{\mathrm{hol}}_{F}\) via the Mellin lift
\(\Upsilon_{F}\) (\Cref{v4-arith:w5-a:thm:prime-side}).
\item The archimedean digamma obstruction as the unabsorbed residual
of the Rankin--Selberg isometry
(\Cref{v4-arith:w5-a:prop:archimedean-not-absorbed}).
\end{itemize}
Chapter~\ref{v4-arith:w5-j:VBstar-spectral-positivity} established:
\begin{itemize}
\item The de Branges structure theorem in the form
\Cref{v4-arith:w5-j:prop:HBplus-iff-kernel-positive}: \(E_{\xi}\) is
Hermite--Biehler if and only if the de Branges reproducing kernel
\(K_{E_{\xi}}\) is positive-semidefinite; with
\((E_{\xi}+E_{\xi}^{*})/2=\Xi\), the Hermite--Biehler separation
theorem puts the zeros of \(\Xi\) on the real axis.
\item The equivalence \(\mathrm{HB}^{+}\Leftrightarrow\mathrm{RH}\)
(\Cref{v4-arith:w5-j:rem:HBplus-alone}), Lagarias 1999.
\end{itemize}
Chapter~\ref{v4-arith:w6-a:ATP-direct-obstruction} established:
\begin{itemize}
\item Unconditional \(W_{\infty}\)-positivity on the high-frequency
dominance sub-class
\(\mathrm{PW}^{\mathrm{hol},\mathrm{HFD}}_{F}\)
(\Cref{v4-arith:w6-a:thm:HFD-ATP}).
\item The bounded-interval form of full A-TP: the negative
archimedean term lies on \([-t_{\ast},t_{\ast}]\), while \(R[f]\)
remains in the inequality
(\Cref{v4-arith:w6-a:thm:bounded-reduction}).
\item Three non-elementary constructions that remain open: zero density,
Gaussian reduction, and de Branges spectral structure
(\Cref{v4-arith:w6-a:rem:live-constructions}).
\end{itemize}

The A-TP\(\leftrightarrow\)de Branges kernel dictionary contains the
Vinogradov--Korobov coordinate change and identifies the obstruction
via Lagarias 1999 and Conrey--Li 2000.

\subsection{Domain}
\label{v4-arith:w7-c:domain}

The objects are:
\begin{itemize}
\item The Lagarias function \(E_{\xi}\), its conjugate
\(E_{\xi}^{*}\), its real part \(A_{\xi}=\Xi\), and the de Branges
kernel \(K_{E_{\xi}}\).
\item The model-space isometry \(K(E_{\xi}^{*}/E_{\xi})\simeq
\mathcal H(E_{\xi})\) and the Mellin transforms which lie in that
model space.
\item Coordinate comparison of the Vinogradov--Korobov zero-free region
under \(s=\tfrac12+iz\).
\item Proof that the Vinogradov--Korobov input gives only peripheral
zero-free bands and no A-TP positivity closure.
\item Identification of the genuine obstruction as HB\(^{+}\),
equivalent to RH.
\end{itemize}

The proof of HB\(^{+}\) on \(\xi_{\mathbb R}\) is RH-equivalent and is
not supplied by the kernel-positivity reformulation. No argument below
uses de Branges operator monotonicity beyond that reformulation.

\subsection{Independent Sources}
\label{v4-arith:w7-c:sources}

The four primary source lists below are disjoint from
the restricted-class comparison and the bounded-interval reduction proof sources.

\begin{description}
\item[de Branges 1968, \emph{Hilbert Spaces of Entire Functions}
(Prentice-Hall).] Foundational monograph. Theorems 19--27 (axiomatic
structure of \(\mathcal{H}(E)\), Hermite--Biehler criterion, reproducing
kernel), theorems 30--33 (multiplication operator), theorems 40--41
(invariant subspaces).
\item[Lagarias 1999, Acta Arith.~89.] \emph{On a positivity property
of the Riemann \(\xi\)-function.} Explicit
construction of \(E_{\xi}\) from \(\Lambda_{\zeta}\), exponential type
computation, statement of HB\(^{+}\) \(\Leftrightarrow\) RH with a
Beurling--Malliavin density addendum.
\item[Conrey--Li 2000, arXiv:math/9812166.] \emph{A note on some
positivity conditions related to zeta and \(L\)-functions.}
Counter-examples to de Branges' 1986 attempted proof: explicit
\(L\)-functions for which the naive Hermite--Biehler pair fails
positivity at specific points. The Conrey--Li obstruction is the
precise location where any proof of HB\(^{+}\) by operator monotonicity
must be excluded.
\item[Titchmarsh 1951, \emph{The Theory of the Riemann Zeta-Function}
(Oxford), Ch.~V--VI.] Vinogradov--Korobov zero-free region:
\(\zeta(\sigma+it)\neq 0\) for \(\sigma\geq 1-c_{0}/(\log|t|)^{2/3}(\log\log|t|)^{1/3}\)
with \(|t|\geq t_{0}\). Korobov 1958 Uspekhi; Vinogradov 1958 Izv.
\end{description}

Independence of inputs. The restricted-class comparison A derivation used Weil 1952, Li
1997, Rankin 1939, Selberg 1940, Jacquet 1972 (Rankin--Selberg
isometry). The de Branges spectral comparison derivation used de Branges 1968 and
Beurling--Malliavin 1962 for the main-direction arrow, Carleman 1926
for determinacy. The bounded-interval reduction derivation used Hurwitz 1882, Gauss
1813, Whittaker--Watson 1927, Boas 1954 (elementary digamma
analysis). The additional sources are de Branges 1968 (in its Ch.~III
reproducing-kernel role, distinct from the structure-theorem role used by the de Branges spectral comparison), Lagarias 1999 (explicit \(E_{\xi}\) computation, not used in the de Branges spectral comparison's structural implication), Conrey--Li 2000 (obstruction location, not
used in the de Branges spectral comparison or the bounded-interval reduction), and Titchmarsh 1951
Vinogradov--Korobov (not used in constructions 5 or 6). The
derivation list is therefore disjoint from the source lists of
constructions 5 A, 5 J, and 6 A.

\section{The de Branges Function \(E_{\xi}\) and Its Real-Line Trace}
\label{v4-arith:w7-c:E-xi}

Fix the shifted character
\(\Xi(z)=\xi_{\mathbb R}(\tfrac{1}{2}+iz)\) of
\Cref{v4-arith:w5-j:def:shifted-character}. We recall the
Hermite--Biehler candidate.

\begin{definition}[Lagarias de Branges function for \(\xi_{\mathbb R}\)]
\label{v4-arith:w7-c:def:E-xi}
\ClaimStatusDefinitional. Define
\begin{equation}
\label{v4-arith:w7-c:eq:E-xi}
E_{\xi}(z)\;:=\;\xi_{\mathbb R}\!\left(\tfrac12-iz\right)
\;+\;\xi_{\mathbb R}'\!\left(\tfrac12-iz\right),
\end{equation}
with its conjugate-partner
\begin{equation}
\label{v4-arith:w7-c:eq:E-star-xi}
E_{\xi}^{*}(z)\;:=\;\overline{E_{\xi}(\overline{z})}\;=\;
\xi_{\mathbb R}\!\left(\tfrac12+iz\right)
\;+\;\xi_{\mathbb R}'\!\left(\tfrac12+iz\right).
\end{equation}
The derivative is with respect to \(s\).  This is the Lagarias
candidate used in the de Branges reformulation.
\end{definition}

\begin{lemma}[Lagarias real-part identity for \(E_{\xi}\)]
\label{v4-arith:w7-c:lem:E-Xi-identity}
\ClaimStatusProvedHere. The real part
\[
A_{\xi}(z):=\frac{E_{\xi}(z)+E_{\xi}^{*}(z)}2
\]
is exactly the shifted xi-function:
\begin{equation}
\label{v4-arith:w7-c:eq:Xi-factor}
A_{\xi}(z)\;=\;\Xi(z)
\;=\;\xi_{\mathbb R}\!\left(\tfrac12+iz\right).
\end{equation}
Consequently \(\mathrm{div}(A_{\xi})=\mathrm{div}(\Xi)\).
\end{lemma}

\begin{proof}
The functional equation gives
\(\xi_{\mathbb R}(\tfrac12-iz)=\xi_{\mathbb R}(\tfrac12+iz)\).
Differentiating \(\xi_{\mathbb R}(s)=\xi_{\mathbb R}(1-s)\) gives
\(\xi_{\mathbb R}'(\tfrac12-iz)=-\xi_{\mathbb R}'(\tfrac12+iz)\).
The derivative terms in \(E_{\xi}+E_{\xi}^{*}\) cancel, and the two
undifferentiated terms give \(2\Xi(z)\).
\end{proof}

\begin{remark}[the completed zeta is not a Hermite--Biehler datum]
\label{v4-arith:w7-c:rem:prefactor}
The function \(\Lambda_{\zeta}(\tfrac12+iz)\) is fixed by
\(*\), by the functional equation.  It cannot satisfy the strict
Hermite--Biehler inequality.  The derivative term in
\eqref{v4-arith:w7-c:eq:E-xi} is therefore not a normalisation; it is
the datum which separates \(A_{\xi}=\Xi\) from its conjugate
quadrature.
\end{remark}

\begin{proposition}[Stirling does not supply Hermite--Biehler positivity]
\label{v4-arith:w7-c:prop:Stirling-asymptotic}
\ClaimStatusProvedHere. No Stirling estimate for the gamma factor in
\(\Lambda_{\zeta}(\tfrac12+iz)\) proves
\[
|E_{\xi}(z)|>|E_{\xi}^{*}(z)|,\qquad \mathrm{Im}\,z>0,
\]
for the Lagarias function \eqref{v4-arith:w7-c:eq:E-xi}.  The omitted
term is the logarithmic derivative of \(\xi_{\mathbb R}\), and the
strict inequality is precisely the Hermite--Biehler condition.
\end{proposition}

\begin{proof}
If \(\Lambda_{\zeta}(\tfrac12+iz)\) were used in place of
\eqref{v4-arith:w7-c:eq:E-xi}, the functional equation would give
\(E=E^{*}\), so the strict Hermite--Biehler inequality would be
impossible.  For the actual Lagarias function,
\[
\frac{E_{\xi}^{*}(z)}{E_{\xi}(z)}
=
\frac{\xi_{\mathbb R}(\tfrac12+iz)+\xi_{\mathbb R}'(\tfrac12+iz)}
{\xi_{\mathbb R}(\tfrac12-iz)+\xi_{\mathbb R}'(\tfrac12-iz)},
\]
and the derivative term is not controlled by the gamma-factor
asymptotic alone.  Thus any proof of the strict inequality is a proof
of Hermite--Biehler positivity, not a consequence of Stirling's
formula.
\end{proof}

\begin{corollary}[no exterior Hermite--Biehler region from Stirling]
\label{v4-arith:w7-c:cor:Stirling-zero-free}
\ClaimStatusProvedHere. The estimate
\[
|E_{\xi}(z)|>|E_{\xi}^{*}(z)|\cdot
\exp\!\left(\tfrac12 r\sin\theta\log(r/2)\right)
\]
is not an unconditional consequence of Stirling asymptotics.  The
valid unconditional zero-free input is the Vinogradov--Korobov
peripheral exclusion for the real part \(A_{\xi}=\Xi\).
\end{corollary}

\begin{proof}
\Cref{v4-arith:w7-c:prop:Stirling-asymptotic} removes the only
purported derivation of the displayed inequality.  The classical
zero-free theorem controls zeros of \(\zeta\), hence zeros of
\(A_{\xi}\), only in the peripheral bands described below.
\end{proof}

\section{de Branges Reproducing Kernel and the Mellin Identification}
\label{v4-arith:w7-c:dB-kernel}

\begin{definition}[de Branges space attached to \(E_{\xi}\)]
\label{v4-arith:w7-c:def:H-E-xi}
\ClaimStatusDefinitional. If \(E_{\xi}\) is Hermite--Biehler, the
\emph{de Branges space}
\(\mathcal{H}(E_{\xi})\) is the Hilbert space of entire functions \(F\)
with
\begin{equation}
\label{v4-arith:w7-c:eq:dB-norm}
\|F\|_{E_{\xi}}^{2}\;:=\;\int_{\mathbb{R}}\left|\frac{F(t)}{E_{\xi}(t)}\right|^{2}\,dt<+\infty
\end{equation}
and \(F/E_{\xi},F^{*}/E_{\xi}\in H^{2}(\mathbb{C}_{+})\), with the
reproducing kernel
\begin{equation}
\label{v4-arith:w7-c:eq:dB-kernel-explicit}
K_{E_{\xi}}(w,z)\;=\;\frac{\overline{E_{\xi}(w)}\,E_{\xi}(z)-\overline{E_{\xi}^{*}(w)}\,E_{\xi}^{*}(z)}{2\pi i(\overline{w}-z)}.
\end{equation}
\end{definition}

\begin{proposition}[de Branges real-line inner product]
\label{v4-arith:w7-c:prop:real-line-inner}
\ClaimStatusProvedHere. For \(F,G\in\mathcal{H}(E_{\xi})\),
\begin{equation}
\label{v4-arith:w7-c:eq:dB-inner}
\langle F,G\rangle_{E_{\xi}}\;=\;\int_{\mathbb{R}}\frac{F(t)\overline{G(t)}}{|E_{\xi}(t)|^{2}}\,dt.
\end{equation}
Here \(E_{\xi}\) is the Lagarias function
\eqref{v4-arith:w7-c:eq:E-xi}, not
\(\Lambda_{\zeta}(\tfrac12+it)\).
\end{proposition}

\begin{proof}
This is the defining inner product of a de Branges space; polarisation
gives the displayed formula.
\end{proof}

\begin{theorem}[Mellin model-space isometry]
\label{v4-arith:w7-c:thm:Mellin-dB-isometry}
\ClaimStatusProvedHere. Assume \(E_{\xi}\) is Hermite--Biehler and put
\[
\Theta_{\xi}(z):=\frac{E_{\xi}^{*}(z)}{E_{\xi}(z)}.
\]
Let
\(\Upsilon_{F}\colon\mathrm{PW}^{\mathrm{hol}}_{F}\to\mathcal{V}^{\mathrm{prim}}_{\mathrm{cusp,temp}}\)
be the Mellin lift of \Cref{v4-arith:w5-a:def:upsilon}. For every
\(f\in\mathrm{PW}^{\mathrm{hol}}_{F}\) whose Mellin transform lies in
the model space \(K(\Theta_{\xi})=H^{2}\ominus\Theta_{\xi}H^{2}\),
define
\begin{equation}
\label{v4-arith:w7-c:eq:dB-lift}
\widehat{\Upsilon}_{F}(f)(z)\;:=\;\widehat{f}(z)\cdot E_{\xi}(z),
\qquad z\in\mathbb{C}.
\end{equation}
Then \(\widehat{\Upsilon}_{F}(f)\in\mathcal H(E_{\xi})\) and
\begin{equation}
\label{v4-arith:w7-c:eq:isometry}
\langle\widehat{\Upsilon}_{F}(f),\widehat{\Upsilon}_{F}(f)\rangle_{E_{\xi}}
\;=\;\int_{\mathbb{R}}|\widehat{f}(t)|^{2}\,dt.
\end{equation}
\end{theorem}

\begin{proof}
For a Hermite--Biehler function, multiplication by \(E_{\xi}\) is the
standard isometric identification
\[
K(\Theta_{\xi})\xrightarrow{\;\sim\;}\mathcal H(E_{\xi}).
\]
Indeed \(F=\widehat f E_{\xi}\) satisfies \(F/E_{\xi}=\widehat f\in
H^{2}\) and \(F^{*}/E_{\xi}=\Theta_{\xi}\widehat f^{*}\in H^{2}\) by
the model-space condition.  By
\Cref{v4-arith:w7-c:prop:real-line-inner},
\[
\|\widehat{\Upsilon}_{F}(f)\|_{E_{\xi}}^{2}=\int_{\mathbb{R}}\frac{|\widehat{f}(t)E_{\xi}(t)|^{2}}{|E_{\xi}(t)|^{2}}\,dt
=\int_{\mathbb{R}}|\widehat{f}(t)|^{2}\,dt,
\]
which is \eqref{v4-arith:w7-c:eq:isometry}.
\end{proof}

\begin{corollary}[A-TP zero form in de Branges coordinates]
\label{v4-arith:w7-c:cor:ATP-as-dB-positivity}
\ClaimStatusProvedHere. Let \(\rho\) run through the non-trivial zeros
of \(\zeta\), and put
\[
z_{\rho}:=-i(\rho-\tfrac12).
\]
Then \(A_{\xi}(z_{\rho})=0\), and \(z_{\rho}\in\mathbb R\) for every
\(\rho\) if and only if RH holds. Under the Mellin--de Branges
model-space isometry \eqref{v4-arith:w7-c:eq:isometry}, the zero-side
A-TP form is represented, after the HB\(^{+}\) realness hypothesis, by
finite Gram sums
\begin{equation}
\label{v4-arith:w7-c:eq:ATP-dB}
\sum_{\rho,\sigma}c_{\rho}\overline{c_{\sigma}}\,
K_{E_{\xi}}(z_{\sigma},z_{\rho}) ,
\end{equation}
with \(c_{\rho}\) obtained from the Paley--Wiener continuation of
\(\widehat f\). Full positivity of the de Branges kernel is the
Hermite--Biehler condition for \(E_{\xi}\); positivity on the sampled
zero set is not a consequence of the change of variables.
\end{corollary}

\begin{proof}
The identity
\[
A_{\xi}(z_{\rho})
=\Xi(z_{\rho})
=\xi_{\mathbb R}(\rho)
=0
\]
is \(\Cref{v4-arith:w7-c:lem:E-Xi-identity}\). If
\(\rho=\beta+it\), then
\[
z_{\rho}=t+i(\tfrac12-\beta).
\]
Thus \(z_{\rho}\) is real for every zero precisely when
\(\beta=\tfrac12\) for every zero. The reproducing property of
\(K_{E_{\xi}}\) writes every finite zero-side sampling form as
\eqref{v4-arith:w7-c:eq:ATP-dB} once the sampled points lie on the
real line. Replacing \(z_{\rho}\) by the real
ordinate \(t\) is therefore legitimate exactly under RH. Finally,
\Cref{v4-arith:w5-j:prop:HBplus-iff-kernel-positive} identifies
full positivity of the de Branges kernel with HB\(^{+}\).
\end{proof}

\section{What the Vinogradov--Korobov Region Actually Gives}
\label{v4-arith:w7-c:VK-strip}

The classical zero-free region is a statement near the line
\(\mathrm{Re}\,s=1\), and by the functional equation near
\(\mathrm{Re}\,s=0\). Under \(s=\tfrac12+iz\) it gives peripheral
horizontal bands in the \(z\)-plane. It does not give a zero-free
strip around the real \(z\)-axis. The real axis is exactly the
critical line.

\begin{definition}[Vinogradov--Korobov peripheral bands]
\label{v4-arith:w7-c:def:VK-strip}
\ClaimStatusDefinitional. Let \(c_{0}>0\) be the Vinogradov--Korobov
constant (Titchmarsh 1951 Ch.~VI; Korobov 1958). Define, for \(T>0\),
the peripheral bands
\begin{equation}
\label{v4-arith:w7-c:eq:VK-strip}
\mathcal{P}_{T}\;:=\;\left\{z=\alpha+i\beta:\ |\alpha|\leq T,\quad
\beta\leq-\frac12+\eta(T)\ \text{or}\ \beta\geq\frac12-\eta(T)\right\},
\qquad
\eta(T)\;:=\;\frac{c_{0}}{(\log T)^{2/3}(\log\log T)^{1/3}}.
\end{equation}
The central strip
\[
\mathcal{C}_{T}:=\{z=\alpha+i\beta:\ |\alpha|\leq T,\ |\beta|<\eta(T)\}
\]
is not a Vinogradov--Korobov zero-free region. It contains the real
axis, hence contains the critical-line zeros when RH holds.
\end{definition}

\begin{lemma}[\(A_{\xi}\) is zero-free in the peripheral bands]
\label{v4-arith:w7-c:lem:E-xi-zero-free-on-strip}
\ClaimStatusProvedHere. For \(T\geq T_{0}\), the function
\(A_{\xi}=\Xi\) has no zeros in \(\mathcal{P}_{T}\). No zero-free
conclusion follows for the central strip \(\mathcal C_T\).
\end{lemma}

\begin{proof}
A zero \(z=\alpha+i\beta\) of \(A_{\xi}=\Xi\) corresponds to a
non-trivial zero
\[
\rho=\frac12+iz=\frac12-\beta+i\alpha
\]
of \(\zeta\). Vinogradov--Korobov gives
\[
\mathrm{Re}\,\rho<1-\eta(|\alpha|)
\]
for \(|\alpha|\) large, hence
\(\beta>-\tfrac12+\eta(|\alpha|)\). Applying the functional equation
gives the reflected inequality
\(\beta<\tfrac12-\eta(|\alpha|)\). Thus zeros with
\(|\alpha|\leq T\) lie in the central band
\[
-\frac12+\eta(T)<\beta<\frac12-\eta(T),
\]
after decreasing \(c_{0}\) and increasing \(T_{0}\) in the usual
uniform way. This proves zero-freeness of \(\mathcal P_T\). It says
nothing about \(|\beta|<\eta(T)\).
\end{proof}

\begin{proposition}[Vinogradov--Korobov gives no A-TP kernel closure]
\label{v4-arith:w7-c:prop:kernel-positive-strip}
\ClaimStatusProvedHere. The zero-freeness of \(\mathcal P_T\) does
not imply positive-semidefiniteness of \(K_{E_{\xi}}\) on the
zero set relevant to A-TP. It only says that the sampled zero set
does not meet \(\mathcal P_T\).
\end{proposition}

\begin{proof}
The A-TP Gram matrix is sampled on the actual zero set, or on its
boundary limits after the Mellin lift. The lemma removes only the
peripheral bands \(\mathcal P_T\). It does not locate zeros on the
real axis, does not put the zero set in a de Branges positivity
domain, and does not identify a positive Hilbert-space inner product
on the sampled zero components. Therefore no A-TP positivity
statement follows from Vinogradov--Korobov alone.
\end{proof}

\begin{theorem}[no Vinogradov--Korobov A-TP closure]
\label{v4-arith:w7-c:thm:ATP-strip}
\ClaimStatusProvedHere. There is no non-vacuous A-TP closure obtained
by restricting Fourier support to the Vinogradov--Korobov peripheral
bands. Any assertion that A-TP closes on a central
Vinogradov--Korobov strip confuses the \(s\)-plane zero-free region
near \(\mathrm{Re}\,s=1\) with the critical-line coordinate
\(\mathrm{Im}\,z=0\).
\end{theorem}

\begin{proof}
If the support lies in \(\mathcal P_T\), the zero sample is empty.
The resulting inequality is the empty inequality and contains no
information about A-TP. If the support lies in the central strip
\(\mathcal C_T\), Vinogradov--Korobov gives no zero-free theorem
there. That central strip contains the real axis, where the
critical-line zeros lie under RH. The proposed closure therefore
has no valid case beyond a vacuous peripheral exclusion.
\end{proof}

\begin{remark}[peripheral support is vacuous for the zero sum]
\label{v4-arith:w7-c:rem:VK-class-nonempty}
Paley--Wiener functions can be arranged to have spectral support away
from the central band, but then the zero-side term controlled by
\eqref{v4-arith:w7-c:eq:ATP-dB} has no zero sample in the
Vinogradov--Korobov region. Such a restriction proves only that a
test which sees no zeros sees no obstruction.
\end{remark}

\section{The Genuine Obstruction: Hermite--Biehler Extension}
\label{v4-arith:w7-c:genuine-obstruction}

\begin{theorem}[Hermite--Biehler obstruction]
\label{v4-arith:w7-c:thm:strip-obstruction}
\ClaimStatusProvedHere. Suppose the Mellin image of
\(\mathrm{PW}^{\mathrm{hol}}_{F}\) is dense in the model space
\(K(E_{\xi}^{*}/E_{\xi})\). Then the extension of
\Cref{v4-arith:w7-c:cor:ATP-as-dB-positivity} from the sampled
zero form to the full de Branges form is equivalent to HB\(^{+}\):
\begin{equation}
\label{v4-arith:w7-c:eq:strip-to-half-plane-equiv}
\bigl[K_{E_{\xi}}\text{ is positive on }\mathcal H(E_{\xi})\bigr]
\;\Leftrightarrow\;
\bigl[E_{\xi}\text{ is Hermite--Biehler}\bigr]
\;\Longrightarrow\;\text{RH}.
\end{equation}
The Vinogradov--Korobov peripheral exclusion is strictly below this
threshold.
\end{theorem}

\begin{proof}
The first equivalence is the de Branges reproducing-kernel criterion
of \Cref{v4-arith:w5-j:prop:HBplus-iff-kernel-positive}.  The
density hypothesis lets positivity on the Mellin image extend to
\(\mathcal H(E_{\xi})\).  Since
\((E_{\xi}+E_{\xi}^{*})/2=\Xi\), Hermite--Biehler separation places
the zeros of \(\Xi\) on \(\mathbb R\), hence gives RH in the
coordinate \(s=\tfrac12+iz\).  The previous section proves that
Vinogradov--Korobov supplies only peripheral zero-free bands; those
bands do not imply the Hermite--Biehler condition and do not touch the
central critical-line problem.
\end{proof}

\begin{corollary}[the genuine obstruction]
\label{v4-arith:w7-c:cor:genuine-obstruction}
\ClaimStatusProvedHere. The irreducible obstruction in the de Branges comparison
formulation of A-TP is HB\(^{+}\) for \(E_{\xi}\), which forces
the Riemann hypothesis through \(A_{\xi}=\Xi\). It is not a lower bound on the
Vinogradov--Korobov width \(\eta(T)\).
\end{corollary}

\begin{proof}
\Cref{v4-arith:w7-c:thm:strip-obstruction}.
\end{proof}

\begin{remark}[comparison to Lagarias 1999 and Conrey--Li 2000]
\label{v4-arith:w7-c:rem:lagarias-conrey-li}
Lagarias 1999 Acta Arith.~89 \S4 proves \(E_{\xi}\)
Hermite--Biehler \(\Leftrightarrow\) RH and discusses the de Branges
space theory of \(\xi_{\mathbb R}\); \S5--\S6 treat the
Beurling--Malliavin density and the analytic structure of
\(\mathcal{H}(E_{\xi})\). The de Branges comparison therefore contributes the
Mellin-to-de Branges placement and the coordinate comparison above: the
Vinogradov--Korobov region is peripheral and cannot close A-TP.

Conrey--Li 2000 arXiv:math/9812166 Thm.~2 exhibits explicit
\(L\)-functions with Siegel-type zeros for which the
\(E_{L}\) candidate does not admit de Branges structure. This confirms
that HB\(^{+}\) on \(E_{\xi}\) is a non-tautological hypothesis: for
an \(L\)-function with a zero off the critical line (analogous to a
Siegel zero), HB\(^{+}\) fails. The Conrey--Li examples identify the
precise location where any proof of HB\(^{+}\) by operator
monotonicity encounters resistance: the absence of such zeros for
\(\zeta\) is the content to be verified.
\end{remark}

\section{Three Non-Elementary Constructions Further Analysed}
\label{v4-arith:w7-c:three-constructions}

Three candidate reductions of HB\(^{+}\) each fall strictly short, by
explicit strength comparison.

\begin{proposition}[Beurling--Malliavin density does not imply HB\(^{+}\)]
\label{v4-arith:w7-c:prop:rule-out-density-only}
\ClaimStatusProvedHere. The Beurling--Malliavin sampling density
argument of \Cref{v4-arith:w5-j:prop:BDplus-reconstruction}
concerns the horizontal spacing of zeros on \(\mathbb{R}\) under RH,
not their confinement to the real axis in the \(z\)-plane.
Beurling--Malliavin density does not imply the Hermite--Biehler
inequality for \(E_{\xi}\).
\end{proposition}

\begin{proof}
BM density is defined on real sequences
(\Cref{v4-arith:w5-j:def:BM-density}); the imaginary-axis
extension uses the projection to real parts, which loses the
vertical information entirely. The horizontal density
\(D_{\mathrm{BM}}(\mathcal{Z}_{\xi})\) can be bounded unconditionally
by the Riemann--von Mangoldt zero count. It does not exclude a single
off-line zero, and hence does not prove HB\(^{+}\).
\end{proof}

\begin{proposition}[Selberg--Bombieri zero density does not imply HB\(^{+}\)]
\label{v4-arith:w7-c:prop:rule-out-polynomial}
\ClaimStatusProvedHere. The Selberg--Bombieri zero-density estimate
\(N(\sigma,T)\leq T^{\theta(\sigma)}\) for \(\sigma\in(1/2,1)\)
with \(\theta(\sigma)\to 0\) as \(\sigma\to 1\) controls the
count of zeros with \(\mathrm{Re}\,\rho\geq\sigma\) in imaginary height
\(T\), but \(\theta(\sigma)>0\) for all \(\sigma<1\). It does not
preclude individual zeros off the critical line. Thus it does not
prove HB\(^{+}\).
\end{proposition}

\begin{proof}
The Selberg zero-density result (Selberg 1946; see Titchmarsh 1951
\S9.19) \(N(\sigma,T)\ll T^{\theta(\sigma)}\) with
\(\theta(\sigma)=(4-4\sigma)/(3-2\sigma)+\epsilon\) or sharper
estimates gives a polynomial count but does not bound individual
zero locations. A density theorem may leave finitely many off-line
zeros, and one such zero already violates HB\(^{+}\). The
density-vs-location mismatch precludes the reduction.
\end{proof}

\begin{proposition}[Connes operator-trace approach reduces to HB\(^{+}\)]
\label{v4-arith:w7-c:prop:rule-out-operator}
\ClaimStatusProvedHere. The Connes adele-class-space construction
(Connes 1999 arXiv:math/9811068; de Branges 1986--1992) has absorption
trace equal to the zero distribution of \(\zeta\); critical-line support
of this determinant-trace datum is RH. Every operator-monotonicity
argument for this support is therefore strictly equivalent to HB\(^{+}\)
and does not constitute an independent reduction below that threshold.
\end{proposition}

\begin{proof}
The Connes trace formula (Connes 1999 Thm.~0.1) identifies the
absorption trace on the adele class space
\(\mathrm{GL}_{1}(\mathbb{A}_{\mathbb{Q}})/\mathbb{Q}^{\times}\)
with the zero distribution of \(\zeta\). Critical-line support of that
distribution is equivalent to RH. Any attempt to prove this support by
operator monotonicity, semi-boundedness, or normality requires an
unconditional proof which is
RH-equivalent, not available at the level of the de Branges comparison. See
\Cref{v4-arith:hpvbconv:rem:connes-form} for the detailed
discussion.
\end{proof}

\begin{remark}[HB ceiling]
\label{v4-arith:w7-c:rem:three-rule-outs}
\Cref{v4-arith:w7-c:prop:rule-out-density-only},
\Cref{v4-arith:w7-c:prop:rule-out-polynomial}, and
\Cref{v4-arith:w7-c:prop:rule-out-operator} establish that no
non-elementary reformulation of A-TP via (i) Beurling--Malliavin
density, (ii) Selberg zero-density estimates, or (iii)
Connes-type operator spectral arguments proves the missing
Hermite--Biehler inequality. The de Branges comparison ceiling is explicit:
Vinogradov--Korobov gives only peripheral zero-free bands, while
HB\(^{+}\) controls the central critical-line problem.
\end{remark}

\section{de Branges Form of A-TP}
\label{v4-arith:w7-c:wave7-form}

\begin{theorem}[de Branges A-TP form]
\label{v4-arith:w7-c:thm:wave7-form}
\ClaimStatusProvedHere. The de Branges form of A-TP is the following.
\begin{enumerate}
\item On the part of \(\mathrm{PW}^{\mathrm{hol}}_{F}\) whose Mellin
transform lies in \(K(E_{\xi}^{*}/E_{\xi})\), A-TP is expressed by
de Branges reproducing-kernel positivity in \(\mathcal H(E_{\xi})\).
\item Vinogradov--Korobov gives only peripheral zero-free bands in
the \(z\)-plane, hence no non-vacuous A-TP closure.
\item A-TP on the full class \(\mathrm{PW}^{\mathrm{hol}}_{F}\) is
bounded by HB\(^{+}\) on \(E_{\xi}\), hence by the RH-strength
realness of \(A_{\xi}=\Xi\).
\item Three non-elementary reformulations (BM density, Selberg
zero-density, Connes operator) do not improve the de Branges comparison
obstruction.
\end{enumerate}
\end{theorem}

\begin{proof}
\Cref{v4-arith:w7-c:thm:Mellin-dB-isometry} gives the model-space
isometry, and \Cref{v4-arith:w7-c:cor:ATP-as-dB-positivity} identifies
the zero-side Gram form.  \Cref{v4-arith:w7-c:prop:kernel-positive-strip}
and \Cref{v4-arith:w7-c:thm:ATP-strip} give the Vinogradov--Korobov
boundary.  \Cref{v4-arith:w7-c:thm:strip-obstruction} identifies the
remaining condition with HB\(^{+}\).
\end{proof}

\begin{remark}[de Branges boundary]
\label{v4-arith:w7-c:rem:consequence}
A-TP is controlled by de Branges reproducing-kernel positivity.
Vinogradov--Korobov supplies only peripheral zero-free bands. The
remaining obstruction is HB\(^{+}\), which is RH-strength through
\(A_{\xi}=\Xi\), and the three
non-elementary reductions stay below that threshold.
\end{remark}

\begin{remark}[relation to the bounded-interval reduction]
\label{v4-arith:w7-c:rem:wave6-relation}
The bounded-interval reduction (\Cref{v4-arith:w6-a:ATP-direct-obstruction}) wrote full A-TP as
\(B[f]\leq A[f]+2\pi R[f]\), with the sign change of the archimedean
kernel confined to \([-t_{\ast},t_{\ast}]\) and
\(t_{\ast}\approx 0.8507\). The de Branges comparison
lifts the archimedean part of this inequality into the de Branges
reproducing-kernel setting. The bounded interval is a window for the
negative sign of \(g\), not the support of the zero set.
Gram-positivity of \(K_{E_{\xi}}\) is imposed on the Mellin image, and
\(R[f]\) remains external. The alleged de Branges sub-range closure does
not exist; the Vinogradov--Korobov bands are peripheral and see no
zero sample. The de Branges comparison therefore does not reproduce
the bounded-interval reduction HFD theorem. It identifies the de Branges form of the remaining
obstruction.

The de Branges comparison refinement over the bounded-interval reduction is analytic: the
bounded-interval archimedean term is transported to a Gram-kernel
inequality on the Mellin image, accessible in principle via de Branges
operator theory. The Vinogradov--Korobov calculation shows that the
usual zero-free region is below the threshold of this problem.
\end{remark}

\begin{remark}[de Branges boundary]
\label{v4-arith:w7-c:rem:honest-form}
The Mellin isometry
\Cref{v4-arith:w7-c:thm:Mellin-dB-isometry}, the de Branges positivity
reformulation \Cref{v4-arith:w7-c:cor:ATP-as-dB-positivity}, and the
Vinogradov--Korobov coordinate comparison
\Cref{v4-arith:w7-c:lem:E-xi-zero-free-on-strip} identify a single
threshold: HB\(^{+}\). Beurling--Malliavin density, Selberg zero
density, and Connes operator spectral reality each remain below that
threshold
(\Cref{v4-arith:w7-c:prop:rule-out-density-only},
\Cref{v4-arith:w7-c:prop:rule-out-polynomial},
\Cref{v4-arith:w7-c:prop:rule-out-operator}). A-TP reduces to
HB\(^{+}\) at the de Branges reproducing-kernel level.
\end{remark}

\section{Consequences}
\label{v4-arith:w7-c:summary}

\begin{enumerate}
\item The de Branges function
\[
E_{\xi}(z)=\xi_{\mathbb R}(\tfrac12-iz)+
\xi_{\mathbb R}'(\tfrac12-iz)
\]
satisfies \((E_{\xi}+E_{\xi}^{*})/2=\Xi\), so the non-trivial zeros
of \(\zeta\) are the zeros of the real part \(A_{\xi}\) in the
coordinate \(s=\tfrac12+iz\)
(\Cref{v4-arith:w7-c:def:E-xi},
\Cref{v4-arith:w7-c:lem:E-Xi-identity}).
\item Stirling asymptotics for the completed zeta do not prove
Hermite--Biehler positivity for \(E_{\xi}\); the strict inequality
is precisely the HB\(^{+}\) input
(\Cref{v4-arith:w7-c:prop:Stirling-asymptotic},
\Cref{v4-arith:w7-c:cor:Stirling-zero-free}).
\item Under HB\(^{+}\), multiplication by \(E_{\xi}\) identifies the
model space \(K(E_{\xi}^{*}/E_{\xi})\) with \(\mathcal H(E_{\xi})\);
the Mellin lift enters this isometry only when its transform lies in
that model space
(\Cref{v4-arith:w7-c:thm:Mellin-dB-isometry}), placing A-TP against
de Branges reproducing-kernel positivity
(\Cref{v4-arith:w7-c:cor:ATP-as-dB-positivity}).
\item The Vinogradov--Korobov zero-free region pulls back to
peripheral horizontal bands \(\mathcal P_T\). It gives no
zero-free central strip around the real \(z\)-axis
(\Cref{v4-arith:w7-c:def:VK-strip},
\Cref{v4-arith:w7-c:lem:E-xi-zero-free-on-strip}).
\item The peripheral zero-free bands do not imply
positive-semidefiniteness of \(K_{E_{\xi}}\) on the zero set. Thus
there is no Vinogradov--Korobov A-TP closure
(\Cref{v4-arith:w7-c:prop:kernel-positive-strip},
\Cref{v4-arith:w7-c:thm:ATP-strip}).
\item The extension to the full \(\mathrm{PW}^{\mathrm{hol}}_{F}\) is
bounded by HB\(^{+}\) on \(E_{\xi}\) and hence by the RH-strength
realness of \(A_{\xi}=\Xi\)
(\Cref{v4-arith:w7-c:thm:strip-obstruction},
\Cref{v4-arith:w7-c:cor:genuine-obstruction}).
\item Three non-elementary reductions (Beurling--Malliavin density,
Selberg zero density, Connes operator) each fall strictly below the
HB\(^{+}\) threshold (\Cref{v4-arith:w7-c:three-constructions}).
\item The bounded-interval archimedean term becomes a
Gram-positivity condition for the de Branges kernel; the
Mellin\(\leftrightarrow\)de Branges isometry of
\Cref{v4-arith:w7-c:thm:Mellin-dB-isometry} is the structural
content underlying both.
\end{enumerate}

\begin{remark}[Beilinson domain boundary]
\label{v4-arith:w7-c:rem:beilinson-closing}
Vinogradov--Korobov 1958 gives a peripheral exclusion after the
change of variables \(s=\tfrac12+iz\). It does not close A-TP.
The obstruction is precise: HB\(^{+}\). Every advance beyond this
requires either a zero-location theorem strong enough to imply
HB\(^{+}\), or an independent proof of HB\(^{+}\) by operator methods.
Both remain open in the arithmetic branch.
\end{remark}

\section{Bibliography Anchors}
\label{v4-arith:w7-c:bibliography}

\begin{description}
\item[de Branges 1968, \emph{Hilbert Spaces of Entire Functions}.]
Prentice-Hall. Ch.~II--III contain the de Branges structure theorem
and the reproducing-kernel formalism used above.
The de Branges comparison argument uses de Branges 1968 in its reproducing-kernel
analytic role (Ch.~III), distinct from the structure-theorem role
used by the de Branges spectral comparison.
\item[Lagarias 1999, Acta Arith.~89.] \emph{On a positivity property
of the Riemann \(\xi\)-function.} Explicit
construction of the \(E_{\xi}\) candidate
(\eqref{v4-arith:w7-c:eq:E-xi}) and Hermite--Biehler formulation
of RH. \S4 is the principal reference.
\item[Conrey--Li 2000, arXiv:math/9812166.] \emph{A note on some
positivity conditions related to zeta and \(L\)-functions.} Explicit
counter-examples where the naive Hermite--Biehler structure fails,
identifying the location where operator-monotonicity arguments
meet an irreducible obstruction.
\item[Titchmarsh 1951, \emph{The Theory of the Riemann
Zeta-Function} (Oxford).] Ch.~V--VI contain the
Vinogradov--Korobov zero-free region used in
\Cref{v4-arith:w7-c:def:VK-strip} and
\Cref{v4-arith:w7-c:lem:E-xi-zero-free-on-strip}.
\item[Korobov 1958, Uspekhi Matematicheskikh Nauk 13.]
Zero-free region \(\sigma\geq 1-c_{0}/(\log|t|)^{2/3}(\log\log|t|)^{1/3}\).
\item[Vinogradov 1958, Izvestiya Akad.~Nauk SSSR Ser.~Mat.~22.]
Independent proof with sharper constants.
\item[Whittaker--Watson 1927, \emph{A Course of Modern Analysis}.]
Cambridge, 4th edition. Stirling asymptotic for \(\Gamma\) used in
\Cref{v4-arith:w7-c:prop:Stirling-asymptotic}.
\item[Connes 1999, arXiv:math/9811068.] \emph{Trace formula in
noncommutative geometry.} Cited for the operator-approach rule-out
of \Cref{v4-arith:w7-c:prop:rule-out-operator}.
\end{description}

The disjointness with constructions 5 A, 5 J, 6 A is per-source. Only
de Branges 1968 appears in both the de Branges spectral comparison and the de Branges comparison; it is used in
distinct roles (structure theorem vs.~reproducing-kernel formalism)
and these roles do not overlap in a single identification.
